\section{Conclusions}
The laboratory demonstrated how an English monoalphabetic substitution can be recovered without exhaustive key search. Frequencies suggested the initial \texttt{e} and \texttt{t} replacements; common patterns such as \texttt{the}, \texttt{that}, and \texttt{to} then constrained the solution. As the passage became readable, longer phrases and specific words resolved the rarer letters.

The 19 recorded steps recovered all 25 symbols present in the 2,567-letter cryptogram. The remaining alphabet pair was inferred by elimination and explicitly distinguished from observed evidence. Each recovered pair was applied consistently throughout the passage.

The central lesson is that frequency analysis is a starting point, not an automatic letter-ranking substitution. Statistical clues become convincing when supported by grammar, word shapes, and repeated occurrences. A large nominal key space does not protect a cipher that retains so much language structure. The exercise also connects with Laboratory Work No. 1: permuting an alphabet may enlarge the key space, but a fixed monoalphabetic substitution still preserves frequencies.

\section*{Report materials}
\addcontentsline{toc}{section}{Report materials}
This laboratory required no application source code. Its materials consist of the ciphertext, recovered key, progress samples, screenshots, and this report, kept under \texttt{Lab\_2} in the laboratory repository:
\begin{center}\url{\repositoryurl}\end{center}
